\chapter*{Conclusions and Further Directions}
\addcontentsline{toc}{chapter}{Conclusions and Further Directions}

This paper has constructed architecture relative to a reading from Atoms and Laws, and has
developed geometry, local consistency, diagnosis, transport, comparison, and reconstruction
as one system. We review the questions that each group of results answered, and state the
mathematical problems that arise next from this system.

\section*{Conclusions}

\begin{itemize}
\item \textbf{Foundational geometry.} We made explicit, as a reading and on the side of the objects, the choice of what to
take as objects and which operations and laws to preserve.
From a finite family of Atoms we generated a core carrying a family of objects closed under
the operations (Theorem~\ref{thm:1.18}), from contexts, covers, and coefficients we
constructed a geometry equipped with a site and a presheaf of rings, and we connected the
three levels of extraction, core, and geometry by projections (Theorem~\ref{thm:1.31}). The
independently defined semantics of lenses and protocols is encoded in this framework, and
the validity of the laws and the semantics-preserving morphisms correspond on both sides
(Propositions~\ref{prop:1.41} and~\ref{prop:1.43}).
\item \textbf{Obstructions to local consistency.} We answered, by an obstruction class, when
states that satisfy the Laws on each part glue into a whole. Satisfaction of the Laws reads
as factorization through the space defined by the ideal of the equations
(Theorem~\ref{thm:2.9}), and when the corrections act as a torsor and the states form a
sheaf, the vanishing of the first obstruction class is equivalent to the existence of a
global state (Theorem~\ref{thm:2.22}). The repair defined on the side of the semantics and
the computation on the side of the equations are connected by a comparison isomorphism of
coefficients, which gives a condition for the existence of a global repair
(Theorem~\ref{thm:2.49}, Corollary~\ref{cor:2.50}).
\item \textbf{Diagnostic invariance.} We gave conditions under which the same obstruction
can be diagnosed even after information is summarized and the ranges read locally are
repartitioned. The canonical resolution that preserves the evaluation of Laws is
characterized by a universal property (Theorem~\ref{thm:3.4}), and the comparison map of
covers, coefficients, and resolutions is an isomorphism on first cohomology under
Condition~C (Theorem~\ref{thm:3.17}). Invariance for all families of Laws for which both
resolutions are adequate can be decided by a finite computation (Theorem~\ref{thm:3.22},
Proposition~\ref{prop:3.23}). Under conditions on the identification of relation components
and on common representatives, the vanishing of the integer-coefficient obstruction and that
of the rational-valued diagnosis agree
(Theorems~\ref{thm:3.37} and~\ref{thm:3.40}). Moreover, for the refinement from three charts
to four charts built from the same Atom input, we constructed the comparison square and
showed that the verdicts, zero or nonzero, agree
(Theorems~\ref{thm:3.43} and~\ref{thm:3.44}).
\item \textbf{Transport and coherence of routes.} We characterized the construction that
carries structure along a change of reading by the universal property that every further
change factors uniquely (Theorems~\ref{thm:4.5} and~\ref{thm:4.11}). The removal of the
discrepancy between a specified comparison and the canonical comparison is equivalent to the
existence of an edge reselection that makes all faces commute at once
(Theorem~\ref{thm:4.23}). For squares generated from finite codes, the exchange of the order
of transport and base change is realized as an invertible canonical comparison
(Theorem~\ref{thm:5.11}), and the range over which one can carry backward is characterized
as the reflection of extraction on the realized locus (Theorem~\ref{thm:5.25}).
\item \textbf{Normalization of comparisons and information about changes.} We factored a
generated comparison involving normalization into an invertible comparison and an idempotent
normalization factor, and showed that invertibility is equivalent to that factor being the
identity (Theorem~\ref{thm:6.16}). Even when it is not the identity, an isomorphism between
the images in the idempotent completion remains (Theorem~\ref{thm:6.12}). For
comparison-preserving changes we gave a necessary and sufficient condition for determination
by observations alone (Theorem~\ref{thm:7.9}), the construction of compatible lifts and the
classification of their totality as a torsor under the kernel
(\S\S\ref{sec:7.4}--\ref{sec:7.6}, Theorems~\ref{thm:7.19} and~\ref{thm:7.21}), and a split
short exact sequence for the operation-preserving changes in systems of operations that keep
hidden states (Theorems~\ref{thm:7.24} and~\ref{thm:7.25}).
\item \textbf{Presentation and local reconstruction.} We showed that, once separation of
morphisms, assembly of morphisms, and assembly of objects are available, the readout functor
gives an equivalence of categories (Theorem~\ref{thm:8.12}), and proved that the category of
full geometries and all their structure-preserving morphisms is equivalent to the category
of local models given by primitive data and independently defined compatibility conditions
(Theorem~\ref{thm:8.18}). The same principle applies to lenses, protocols, and design
descriptions divided among teams, and gives the unique extension from finite tables
(Propositions~\ref{prop:8.26} and~\ref{prop:8.28}) and the distinction between information,
procedure, and computational cost for finite models (\S\ref{sec:8.8}).
\end{itemize}

The questions of the introduction, the freedom of change and the information needed for a
judgment, received the following answers.
The degrees of freedom of the permitted changes are determined by how far the operations
tie the correspondences between states together, and for systems of operations that keep
hidden states the number of such changes is given by the families of permutations indexed
by connected components together with the torsor under the kernel (\S\S\ref{sec:7.6}--\ref{sec:7.8}). The information needed for a
judgment is characterized by the inclusion of the kernel of the observation
homomorphism in the comparison-preserving group (Theorem~\ref{thm:7.9}).
Even when compatible lifts can be constructed from the normalized information, judging a
given change may require information that normalization has lost
(\S\S\ref{sec:7.4}--\ref{sec:7.5}).

These results were proved starting from constructions on the same input and using the output
of an earlier result as the input of the next. The answers to the questions of local consistency, diagnosis,
transport, comparison, and reconstruction are connected to one another as theorems about
common objects.

\section*{Further directions}

From the objects and constructions of this paper, the following problems can be formulated.

\begin{enumerate}
\item \textbf{Coherence of reconstruction and comparison groups.}
Chapters~\ref{chap:6}--\ref{chap:7} constructed the normalization of generated comparisons,
the group of comparison-preserving changes, and a section of the restriction homomorphism.
Chapter~\ref{chap:8} gave the recovery of objects and all structure-preserving morphisms.
The next question is how to give, for an arbitrary generated comparison, the
correspondence
between its comparison-preserving group and the comparison-preserving group after
normalization in a form naturally compatible with the projections to the base, the
observations, and the coefficients, and to derive it all at once from the recovery given by
the main equivalence (Theorem~\ref{thm:8.18}), including compatibility with the subgroup
that fixes the base, with idempotent completion, and with the category of morphisms. The
sections, split short exact sequences, and descriptions of kernels and fibers constructed
individually in Chapter~\ref{chap:7} become checkpoints for this general correspondence.
\item \textbf{A common definition of distinction, extension, and effectiveness.} In
\S\ref{sec:8.8} we distinguished the information that determines objects and morphisms, the
checking procedure based on enumeration and equality tests, and the computational cost. The
next task is to give a common definition for these and to establish it as a necessary and
sufficient condition in terms of the connected components of the system of operations
(\S\ref{sec:7.7}) and as a computation of determination and extension over explicit finite
inputs.
\item \textbf{Unification of finite determinacy.} Proposition~\ref{prop:8.32} showed the
limits of a countable syntax with respect to automorphisms of the extraction base.
Proposition~\ref{prop:8.33} showed that, even when a family of tags is recovered from all
finite fragments, it is not determined by a single finite fragment over an infinite index.
For lenses and protocols we obtained the unique extension from finite tables
(Propositions~\ref{prop:8.26} and~\ref{prop:8.28}) and the connection with retracts,
idempotent completion, and comparison squares (Proposition~\ref{prop:8.31}). The question is
to unify these as recovery by the main equivalence under the same readout, and to state, as
a single theory of finite determinacy, the judgment of invertible changes by connected
components together with the recovery of comparison groups, sections, kernels, and fibers.
\item \textbf{Descent for families and global representation.} The gluing of
Chapter~\ref{chap:2} and the obstruction to composition of Chapter~\ref{chap:4} treated
obstruction classes for each fixed input. Once sites, coefficients, and transport are
indexed by families, one can ask for the relation between the vanishing of the obstruction
for each family and globalization, the correspondence between the first cohomology that
classifies lifts of local changes and the existing obstruction classes, and Morita-type
equivalences between categories of comparisons. Beyond that we place the problem of
representability: whether the functor of realizations, natural in the coefficient algebra,
is represented by a geometric object such as a scheme or a stack.
If representability holds, the totality of realizations for a fixed input becomes a single
geometric object, individual realizations can be treated as its points, and families of
realizations as morphisms into that object. The range of repairable realizations, and the places
where the value of the diagnosis jumps, can then be stated globally as computations of
subsets of that object.
These objects and assumptions are formulated separately from the problems treated by the
foundational constructions of this paper.
\item \textbf{Structure in the time direction.} The reading used in this paper concerns the
structure of a single version. Introducing traces subordinate to a sequence of measurements,
a product site, and coefficients in the time direction, one can ask under which assumptions
transport and gluing along time can be compared with the constructions of transport and
descent of this paper.
\item \textbf{Connection with measurement and with a theory of evolution.} The input of this
paper consists of observed Atoms, a specification of Laws and coefficients, and a fixed
reading. ArchSig, the author's measurement tool, computes diagnoses from two streams of
input: observation records of source code, and a specification of laws and equations. To
connect this measurement with the present paper, one must formulate as verifiable conditions
that the extracted observations satisfy the conditions on extraction of
Chapter~\ref{chap:1}, and that the computed diagnoses satisfy the assumptions on comparison
of Chapter~\ref{chap:3}. Once this connection is complete, the judgments of each chapter of
this paper become measurements with evidence, returning the same result from the same input
on actual source code. Besides the detection of inconsistencies that are correct locally but
break as a whole (Chapter~\ref{chap:2}) and the sharing of that verdict along a replacement
of the unit of review (Chapter~\ref{chap:3}), the scope of measurement extends to the
computation of the discrepancies between several migration routes and the decision whether a
reselection that removes them exists (Chapter~\ref{chap:4}); the check that the procedure of
splitting the whole and then selecting a scope agrees with the procedure of selecting a
scope and then splitting (Chapter~\ref{chap:5}); the decision of what a normalization to the
standard representation preserves and whether one can return to the original
(Chapter~\ref{chap:6}); the determination by observations whether a proposed change
preserves a comparison, and the number of all compatible candidate changes
(Chapter~\ref{chap:7}); and the finite check whether design descriptions divided among teams
assemble into a single model and a structure-preserving morphism (Chapter~\ref{chap:8}).
For this last check, Chapter~\ref{chap:8} even distinguishes the information
required, the procedure, and the computational cost.
Furthermore, for the connection with a theory that treats how artifacts and practices change
the reachable future of software evolution (Software Field Theory), the question is to
formulate a correspondence that takes the freedom given by the classification of changes of
Chapter~\ref{chap:7} and by the reconstruction of Chapter~\ref{chap:8} as the input of a
description of evolution. Once this correspondence is obtained, the process of development
is described as motion in the space formed by all realizations, and one can ask, as an
object of computation, which changes and practices widen the range of reachable designs and
which narrow it.
\end{enumerate}

Every one of these problems can be stated from the objects and maps constructed in this
paper. The Rising Sea approach, of enveloping individual problems by raising the water level
of the theory, continues on this foundation. The higher the water rises, the more the problems
of development practice acquire coordinates, become theorems, and descend to measurement.
